\documentclass[man,11pt,a4paper]{apa7}
% APA 7 professional-manuscript layout, set in a sans-serif (Helvetica, metrically
% identical to Arial). Converted from the Word original; wording unchanged.
\usepackage[T1]{fontenc}
\usepackage[utf8]{inputenc}
\usepackage[scaled]{helvet}
\renewcommand{\familydefault}{\sfdefault}
\usepackage{enumitem}
\setlist[enumerate]{itemsep=0pt, topsep=0pt}

\title{Risk, Relevance, and the Aggregation of Competing Claims}
\shorttitle{RISK, RELEVANCE, AND COMPETING CLAIMS}
\authorsnames{SJ Beard}
\authorsaffiliations{{Centre for the Study of Existential Risk, University of Cambridge}}
\authornote{This paper was posted in 2026 but is based on material written in 2018 and 2019 and has not been updated since. Subsequent work that built on these ideas was published in other papers, including Beard, S., \& Kaczmarek, P. (2019). On the wrongness of human extinction. \textit{Argumenta}, \textit{5}, 85--97.. This work was facilitated by the Centre for Philosophy of Natural and Social Science at the London School of Economics, who hosted the author as a visitor in 2018.}
\abstract{This paper seeks to expand on Alex Voorhoeve's Aggregate Relevant Claims view to enable it to deal with cases involving known levels of risk to actually existing people. Voorhoeve's view involves two concepts, the strength of a claim and its relevance relative to other claims, which he argues are closely connected. In the first section of this paper, I examine how the concepts of strength and relevance might interact with risk and show that, if interpreted crudely, his view would face significant problems in such cases. I argue that we can make the Aggregate Relevant Claims view perform much better in relation to problems of risk by 1) separating the notions of relevance and strength and 2) redefining relevance not as a single number, analogous to strength, but as a range concept with upper and lower bounds of relevance that can move independently. Next, I argue that the same move can be used to improve upon the Aggregate Relevant Claims View's approach to assessing the relevance of claims whose strength arises from inequalities. Finally, I highlight some general conclusions about how my arguments can produce radically different approaches to aggregating competing claims under different conditions of risk and inequality, and why this seems to me intuitively attractive. In an appendix, I discuss whether the view can also be applied in terms of existential risk, in which the claims of possible future people are of particular concern.}

\begin{document}
\maketitle
% SJ's choice (not APA): single-spaced footnotes
\makeatletter
\renewcommand\@makefntext[1]{\linespread{1}\selectfont\raggedright\textsuperscript{\@thefnmark}~#1}
\makeatother

\section{Risk and the Aggregate Relevant Claims View}

According to `Competing Claims' views about fairness, we should evaluate alternatives based on the comparative strength of the claims of different individuals. These views were designed to address issues in distributive justice, and in particular the `Separateness of Persons' whereby we seem to value well-being differently when making interpersonal tradeoffs between people and making intrapersonal tradeoffs between potential futures for one person (Voorhoeve \& Fleurbaey, 2012). This means that they tend to treat issues of fairness and risk separately; however, in real world cases they often co-occur.

This paper considers the application of one such theory, Aggregate Relevant Claims (ARC) developed by Alex Voorhoeve, to cases involving both unfairness and risk. This view was developed to mediate between, and allow us to quantify, both pro and anti-aggregationist leanings concerning the evaluation of outcomes, which could be especially important in evaluating certain pressing global challenges, because the people who need to coordinate their behaviour to address them often span groups traditionally opposed to the aggregation of individual claims (such as public health professionals and lawyers) and those who support it (such as economists and engineers). Whilst it is possible to engage with these groups separately, developing bespoke moral arguments for each, the need to develop a rational and coherent response to these challenges gives us reasons to want to develop moral theories that can speak to both groups, and that can help us to understand when their disagreements are most important and when they are not.

ARC is expressed in the following five axioms:

\begin{enumerate}[label=\arabic*.]
\item Each individual whose well-being is at stake has a claim on you to be helped. An individual for whom nothing is at stake does not have a claim.
\item Individuals' claims compete just in case they cannot be jointly satisfied.
\item An individual's claim is stronger:
\begin{enumerate}[label=\alph*.]
\item the more her well-being would be increased by being aided and;
\item the lower the level of well-being from which this increase would take place.
\end{enumerate}
\item A claim is relevant if and only if it is sufficiently strong relative to the strongest competing claim.
\item You should choose an alternative that satisfies the greatest sum of strength-weighted, relevant claims.
\end{enumerate}

Note that whilst axiom 3 states the direction of the relationship between the strength of a claim and facts about the well-being of the individual who possesses it, it does not indicate other features about this relationship. If we assume that the relationship is non-linear and increasing (as I will tend to do) then the theory will adopt a prioritarian attitude, showing greater concern for the satisfaction of stronger claims over weaker ones. If it is assumed to be linear then the view will still be prioritarian about the distribution of well-being, because it still shows more concern for the well-being of the worst off, but to a lesser extent and without giving priority to the satisfaction of stronger claims over weaker ones. If we assume that the relationship is non-linear and decreasing then the theory will give priority to satisfying weaker claims (where these are relevant), though that would seem somewhat strange given the other features of this view.

As with many views about distributive fairness, ARC was developed primarily with reference to the distribution of whatever we take to be valuable across people (commonly labelled `well-being') and was specifically developed for cases where the distribution of well-being was known for certain (Voorhoeve, 2014, p. 67). This paper considers how to expand the ARC to deal with cases involving risk. In particular, it deals with the issue of interpersonal (or social) risk. For instance, how should we aggregate the claims of individuals who are certain to suffer from a particular illness or other harm and those who merely face the risk of doing so?

In order to apply the ARC view to such cases, we need to determine how the strength of a person's claims, and their relevance to other claims, depends upon the level of risk they face. According to the view the strength of a claim is determined by two things, the size of benefit the person might receive and their well-being level, while the relevance of the claim is determined by its strength. When a benefit or cost is uncertain, however, should we understand its size to depend on its probability, for instance the probability that they will suffer from some avoidable risk? To put this another way: If we save someone from a 1 in 10 chance of death, should we treat this the same as if we had saved their life, or as if we had given them a lesser benefit?\footnote{In all the cases I consider, I will describe the probability of a risk and benefit ex-ante, which I take to be the point of view of the decision maker.}

\section{Adapting the ARC}

The first possible way in which ARC could evaluate this kind of risk would be if the size of benefit that a person receives were to depend upon how much they might benefit (ex-post), but not the probability that they would benefit at all (ex-ante). However, this seems to provide counterintuitive results with regard to which claims are the strongest. For instance, if I can benefit one person by some amount, x, for certain, or instead have a one in 100 chance of benefiting two people by the same amount, it seems wrong to say that I should (or even could permissibly) act so as to have a slim chance of benefiting the two people because the combined strength of their claims is 2x rather than merely x. I take the force of this objection to be obvious -- the probability that a person will benefit must discount the size of their claim to a benefit in some way. However we still need to determine what this way is.

The second way in which the ARC could respond to risk would be if the size of benefit that a person receives depends on the size of their expected benefit, i.e. it took equal account of both how much they would benefit and their probability of benefiting. This may also have counterintuitive implications, albeit this time with respect to the relevance of claims, which according to the ARC is a simple function of their strength. Let us suppose that in the previous example I now have to choose between certainly benefiting 1 person by x or having a 1 in 100 chance of benefiting 1 million people by x. Let us also assume that, under the ARC, if one claim is only 1/100\textsuperscript{th} the strength of another claim then it is not sufficiently strong to be relevant when the two claims compete. Now we are left with the conclusion that, whilst the combined benefit of helping the million people is expected to be 10,000 times greater than that of helping the 1 person -- and that we can expect to provide 10,000 people with the same sized benefit in this case (albeit we may in-fact benefit more or less than this but with a vanishingly small chance of benefiting nobody); we should still act to certainly benefit the 1 person, because none of the claims of the 1 million is sufficiently strong to be relevant, so their combined strength is irrelevant for the purposes of aggregation.

However, this implication does not force us to the conclusion that providing 10,000 times the benefit by helping the 1 million people means satisfying claims that are collectively 10,000 times as strong. As I suggested earlier, when considering the strength of competing claims we can apply some priority weighting on the strength of claims, such that a claim to receive double the benefit may be more than twice as strong. We might therefore perfectly well claim that providing 100 people with a 1 in 100 chance of benefiting by x would satisfy claims that were nowhere near as strong as providing 1 person with a certain benefit of x. Indeed, I shall assume for the rest of this paper that we hold such a view. However, it is likely that, in this case, providing a benefit of 10,000x would mean satisfying claims that were, collectively, at least as strong as the single claim of one person to have a certainty of benefiting by x.

A similar example, albeit involving a risk that is purely epistemic in nature, has been considered by Michael Otsuka who comes to the same conclusions. Otsuka's case involves a policy choice between sacrificing the limb of a known Floridian and exposing every single Californian to a risk of death such that we know that exactly one Californian will die, but not who that person will be. Otsuka argues that in such cases we should not discount the claim of each Californian in accordance with the miniscule (1 in 40,000,000) risk of their own death because:

\begin{quote}
Once it has been so heavily discounted, each Californian's complaint against death becomes something fairly trivial: the complaint against a subjection to a miniscule increase in one's risk of premature death that is no greater than various risks of premature death we accept as a matter of course in exchange for small conveniences. Such complaints might be regarded as so trivial as to be morally irrelevant when compared with serious complaints. (Otsuka, 2015, p. 80)
\end{quote}

I find Otsuka's reasoning compelling in this case and believe it to apply to non-epistemic kinds of risk as well. There seems to me to be a clear difference between a claim that is small because it is a claim to a very unlikely but significant benefit and a claim that is small because it is a certain claim to a trivial benefit such as forgoing some annoyance, such as T. M. Scanlon's (1998) example of an interrupted World Cup game. Whilst we need a moral theory that discounts the size of a claim according to its probability of occurring we also need a theory that recognizes this important distinction.

In searching for such a theory Otsuka suggests the following

\begin{quote}
In my case of Californians at risk, the trivially low risks of death shared by millions together constitute something very far from trivial: the certainty of a death. By contrast, millions of viewers' experiencing annoyance and frustration over an interrupted World Cup match do not together constitute any serious harm to anyone. But this difference... reveals that it is superfluous to aggregate the individually minor complaints of the many Californians. Rather, it is sufficient merely to appeal to the strength of the nondiscounted complaint of the one Californian who would die. (Otsuka, 2015, p. 80)
\end{quote}

However, I disagree. I do not think that discounting risks according to the probability that `someone' will be harmed rather than the probability that each person will be harmed is correct for a number of reasons. The first is that, as Otsuka admits, people do accept miniscule increases in the probability of premature death in exchange for miniscule benefits as a matter of course (for instance when crossing the road). However, Otsuka's view implies the opposite, if the government were in a position to prohibit people from doing this then they should, because the claim of the small number of people who actually die as a result of these exchanges would render all the small benefits people seek by accepting such a risk trivial.

My second objection to Otsuka's view is that it seems to suggest that exposing 100 people to 1 in 100 risk of suffering a certain harm must lead to claims that are collectively about as strong as exposing one person to a certainty of suffering that harm, because that is what we can expect to be the eventual consequence of our action. However, whilst I agree that we should not view a small risk of a great harm as trivial when compared with a certainty of a lesser, but still significant harm, the claim that the collective claims of 100 people who each face a 1 in 100 risk of being harmed must collectively be equally strong as that of a single person facing a certain harm appears to me to be too strong. As I suggest above, the ARC could give us a way of preserving this intuition by prioritizing the satisfaction of stronger claims over weaker ones, at least to some extent.

\section{Modifying the Notion of a Claim's Relevance}

Together, these cases suggest that the strength and relevance of a person's claims should be separated in cases of social risk, or low probability costs will either generate claims that are too strong or insufficiently relevant relative to high probability costs. One way of doing this might be simply to modify ARC to make the relevance of a claim depend on both the size of the benefit and the well-being level of the person being benefited but not the probability that the person will benefit, even though the strength of a claim depends on all three. On this revised view, any claim to receive a benefit of X would be equally relevant irrespective of how likely one is to be benefited. However, the strength of the claim, and hence the ultimate determinant of what we should do, would still depend on the size of the benefit, the well-being level of the person to be benefited and the likelihood that they will be benefited, so that claims to certain benefits will still be a lot stronger than claims to merely probable or even unlikely benefits.

Such a view would imply both that we should give 1 person a certainty of benefiting by x, rather than 2 people a 1 in 100 chance of benefiting by x, because the claim of that 1 person is much stronger than the combined claims of the two, and that we should give 1 person a certainty of benefiting by x rather than 100 people a 1 in 100 chance of benefiting by x, because whilst the size of the two expected benefits would be equal, the claims of the 100 would be weaker. However, it would not imply that we should give 1 person a certainty of benefiting by x rather than a million people a 1 in 100 chance of benefiting by x, because the claims of the million are still relevant, and even though the strength of their individual claims may be much smaller the total strength of their combined claims may be much greater. However, it thus respects Otsuka's intuitions that one should not simply dismiss as irrelevant low probability but significant benefits, in the same way as one might certain but trivial ones.

However, it still would not be a sufficient modification of the ARC in order to avoid problematic implications when dealing with cases involving risk. This is because it only works for cases involving the relevance and strength of claims to be benefited by equal amounts (amount x) but works much less well in situations where the amount one might benefit by differs.

For instance, let us assume that losing a toe is not a relevant claim compared to death, such that we would not believe that we should save any number of people from losing a single toe if we could instead save just one person from death. If the relevance of a claim were to depend only on the amount that somebody stood to benefit, and not the probability of their benefiting, then it would seem that we have to say that losing a toe is also not a relevant claim when compared to a 1 in 100 chance of death, or even a 1 in 1,000,000 chance of death, even though such claims may be weaker than a claim to lose a toe. Indeed, as we have already discussed in relation to Otsuka's view, individuals accept such trade-offs as a matter of course. Parents, for instance are willing to accept surgery for their children to correct minor defects and ailments, even though even comparatively safe surgeries carry a not insignificant risk of death from anaesthetic, deep vein thrombosis and so on, if for no other reason than because the expected benefits of such surgery far outweigh these expected costs.

A further problem is that it seems intuitive that, for any theory about the aggregation of claims, a stronger claim should always be at least relevant when compared to a weaker claim. Yet if we were to totally divorce the concept of a claim's relevance from its strength, then this can no longer be guaranteed as cases such as these appear to show. I find this troubling.

There may be several solutions to this problem. However, the one I wish to explore in this paper is that we should abandon the notion that a claim's relevance refers to any single value, whether directly related to that claim's strength or not. Instead, we should redefine a claim's relevance as a range. According to the ARC, as modified above, every claim has a single relevance value, and what we care about is the 'relative relevance' of any two competing claims. However, other views might instead take relevance as a range concept, such that each claim's relevance is determined by two numbers, the strongest claim compared to which it is still relevant (and above which it becomes irrelevant when compared with all other claims) and the weakest claim compared to which other claims are still relevant compared to it (and below which all other claims become irrelevant when compared to it).\footnote{Let us assume that two claims are relevant to one another if their relevance ranges overlap. There is an alternative view we might take in which two claims are relevant if the strength of one claim falls into the relevance range of the other. I prefer the interpretation that I give here because it is more symmetrical and because it allows that two claims that have very different strengths may be relevant to one another under conditions that increase their relevance ranges. Others may take this view as too permissive of aggregation, and prefer the alternative interpretation. At present, I see no strong arguments to judge between them.}

This modified view could be stated as follows:

\begin{enumerate}[label=\arabic*.]
\item Each individual whose well-being is at stake has a claim on you to be helped. An individual for whom nothing is at stake does not have a claim.
\item Individuals' claims compete just in case they cannot be jointly satisfied.
\item An individual's claim is stronger:
\begin{enumerate}[label=\alph*.]
\item the more her well-being would be increased by being aided;
\item the more likely her well-being would be increased by being aided and;
\item the lower the level of well-being from which this increase would take place.
\end{enumerate}
\item A claim's relevance consists of a range:
\begin{enumerate}[label=\alph*.]
\item whose upper bound is determined solely by the amount of well-being increase a person would receive from being aided and the well-being level of that person and
\item whose lower bound is determined by the amount of well-being increase that person would receive, the well-being level of that person \textit{and the probability that their well-being would be increased}.
\end{enumerate}
\par a claim is relevant if its relevancy range overlaps with the relevance range of the strongest competing claim.
\item You should choose an alternative that satisfies the greatest sum of strength-weighted, relevant claims.
\end{enumerate}

If we understand relevance in this way, then we no longer have to say that probability does not affect the relevance of a claim, all we have to say is that probability does not alter the upper bound for the strength of other claims with which it is relevant. However, the lower bound for which claims would be relevant to compete against this claim can still decrease as a benefit becomes less likely and hence the strength of the claim reduces. Another way of putting this is that when a claim is weakened because it involves an unlikely benefit the claim will remain relevant when compared with the same strengths of claim that it would have been even if it involved a certain benefit, but many more claims will now be relevant when compared with it to the extent that any claim that would be relevant when compared with a claim to a certain benefit of the same strength as this claim will also be relevant when compared with this claim. Note that on this view the strength of a claim will always fall within its relevance range, thus preventing the result that a claim may not be relevant even when compared with another claim that is weaker.

On this view then, we might say that even though losing a toe is not a relevant claim when compared with a certainty of death, a 1 in 100 chance of death is both relevant when compared to a certainty of death because they involve benefits of the same size, even though one is much less likely (and hence much weaker), and also that losing a toe is relevant when compared with a 1 in 100 chance of death because they involve claims of comparable strength, even though one involves a benefit of a much greater size. Indeed, on my preferred interpretation of the view we would go so far as to say that a 1/100 risk of losing a toe would be relevant when compared with a 1/100 risk of death, because even though it involves a claim that is much weaker the relevance ranges of the two claims still overlap.

This seems to me like the best set of solutions to these problems involving social risk, and this modification to the ARC view seems able to provide interesting new avenues for considering important issues in the ethics of risk, such as the relative values of identified versus statistical lives.\footnote{A different problem concerns how to evaluate cases involving deep uncertainty (also known as ambiguity), where even the probability of outcomes cannot be known. Such cases are often modelled in terms of uncertainty about the underlying `state of nature' and I propose that we might plausibly handle them by treating people's claims across possible states of nature as if they were the claims of the different people in each of these states (which might compete against each other). One reason to think this is the Rawlsian view that in situations of deep uncertainty people should act, in their own self-interest, as if they were certain to face the worst possible outcome rather than weighing costs and benefits across outcomes as one normally does. Another is the principle of uncertainty aversion, whereby decision-makers are morally obliged to take extra precautions against the worst possible implications of their actions, even if this means the same people miss out on substantial benefits. If we allow that individuals' claims may compete both interpersonally and intrapersonally across different states of nature, while within any one state of nature individuals' claims compete only interpersonally but not intrapersonally (for instance, by changing the ARC's first axiom to read `Each individual in a state of nature whose well-being is at stake...'), this opens up further, potentially productive, ways ARC could deal with more complex, real world problems.}

\section{Can This Modified ARC Meet Seth Lazar's Desiderata for Aggregation Under Risk?}

Seth Lazar (2018) sets out four desiderata for any view that aims to contain anti-aggregationist elements to meet when applied to cases involving risk. These are as follows

\begin{enumerate}[label=\arabic*.]
\item When the risked harms are relevant to one another, lower individual risks of harm may be aggregated to outweigh higher individual risks of harm.
\item When the risked harms are not relevant to one another, and the probability of the greater harm is high enough, individual risks of lesser harms may not be aggregated to outweigh the risk of the greater harm.
\item When the risked harms are not relevant to one another, and the probability of the greater harm is low enough, individual risks of lesser harms may outweigh a small enough risk of the greater harm.
\item When facing a series of choices to spare people from a risk of a lesser harm, each of which imposes a very low risk of a greater harm, but where the series has a very high risk of the greater harm befalling someone, the individual risks of lesser harms may be aggregated to outweigh the cumulative risk of the greater harm. (Lazar, 2018)
\end{enumerate}

The modified version of the Aggregate Relevant Claims view that I have set out above is able to meet all of these desiderata. Here is how it would be applied in the cases suggested by Lazar.

For instance Lazar offers the following case as an illustration of desideratum 1: If you $\varphi$, then one person will face a 0.5 probability of death, whilst 10 million people will be fine; if you $\psi$, then these 10 million people face a 1 in 1 million risk of death whilst the original person will be fine. In this case, desideratum 1 requires that we may have reason to save the 10 million from their very slight risk of death, even though the strength of their individual claims are so much weaker than that of the single person facing a 0.5 chance of death. This is indeed what the theory I am proposing would imply, though whether or not it would say that in this case there are enough of the lower individual risks to fully outweigh the higher individual risk is undetermined, it will depend upon the priority weighting we apply to satisfying strong claims rather than weaker ones. However it will certainly be the case that there is some number of lower individual risks such that people's claims to avoid these risks can aggregate to outweigh that of a single person's claim to avoid a higher risk. The relevance of the claims to avoid these risks is ensured, on this theory, because they are all risks to avoid the same negative outcome, namely death.

The illustration of desideratum 2 that Lazar offers is as follows: Suppose that you have to choose between leaving one person exposed to a 0.9 risk of death whilst fully curing a very large number of other people; or fully cure this one individual of their risk of death whilst leaving everyone else exposed to a 0.9 risk of a minor headache. In this case, desideratum 2 requires us to save the one person no matter how many people we could save from the headache. Again this is what the theory I propose here will recommend. Whilst the claim to be cured of a 0.9 risk of death may be relevant to a slightly larger set of claims than that to be cured of a certainty of death, since the risk of death is very high then the strength of the claims to which this will be relevant will only extend to include a few slightly weaker claims, and certainly not to claims that involve merely a risk of headache. Therefore, this theory will claim that since the claim to be cured from a 0.9 risk of a headache is not relevant when compared with the claim to be saved from a 0.9 risk of death then we cannot aggregate the weaker claims in order to outweigh this one stronger claim.

For Desideratum 3 Lazar asks us to consider having to choose between 1 individual facing a 1 in 10,000 risk of death, whilst a very large number of people will be fine and 1 person being fine whilst a very large number of people will face a 0.9 risk of a minor headache. Here, desideratum 3 implies that we should leave the one individual to face their 1 in 10,000 risk of death. It is plausible that the theory I am proposing here will have the same implication, since the very low risk of death implies that a person's claim to be saved from this risk, whilst still relevant to claims to be saved from a much higher risk of death, will also be relevant when compared with claims to be saved from much smaller risks, including, plausibly, risks of minor headaches. If the two claims are relevant to one another then, so long as there are enough people with a claim to be saved from the risk of a headache, regardless of its size, they can outweigh this one claim to be saved from a very low risk of death. (Lazar, 2018)

Finally the case Lazar uses to illustrate desideratum 4 involves merely repeating the case presented for desideratum 3 multiple times. Desideratum 4 implies that this will not change the result of this case, even though the risk that somebody will die is now very close to 1. However such repetition will have no effect on the implications of the theory I have proposed here since the relevance of claims depends upon how they affect individuals, rather than groups.

Lazar proposes a different approach to aggregation under risk that is also able to satisfy these desiderata. At present, it is not clear how similar our two views are. However, they share the common feature that, in order to assess the relevance of claims, we need to know facts about both their ex-ante and ex-post effects. Lazar's main objection to his own view however is that he believes it may be `inelegant' (Lazar, 2018). Such things are of course highly subjective, however my own sense is that the modified Aggregate Relevant Claims view that I present here may perform slightly better on that metric.

\section{How the Modified ARC Evaluates Inequality}

Our modified ARC view thus provides a good theory for evaluating cases involving social risk. Let us now return to the kind of cases for which the ARC was originally developed, and consider whether these modifications might also be useful for such cases as well.

Can one claim render another claim of the same, or similar, strength irrelevant merely because of an inequality between the well-being levels of the people being aided? Alex Voorhoeve has said (in personal correspondence) that he does not believe that this is possible and this appears like the correct conclusion to me. However, it seems hard to arrive at such a conclusion under the ARC, as defined by Voorhoeve, because on this view inequalities do have a considerable impact on the strength of claims and this is supposed to translate directly into their relevance.

For instance, it seems like many people believe that when comparing the following two populations

\begin{tabbing}
Population A:\quad\= \kill
Population A: \> 1 billion people at level 100, 200 billion people at level 1 \\
Population B: \> 201 billion people at level 2
\end{tabbing}

\noindent The fact that population B is more equal than population A is a good thing, even though the benefits to the worst off in A compared to B are so much smaller than the losses of the best off in A relative to B, as to render them irrelevant when compared with the best off people in A compared with B. If, on the ARC, the fact that these benefits come to people is supposedly of such moral significance as to make these otherwise irrelevant claims relevant then it seems like they should be sufficiently morally significant to make otherwise relevant claims of the better off irrelevant. For instance consider this alternative case

\begin{tabbing}
Population A:\quad\= \kill
Population C: \> 200 billion people at level 100, 1 billion people at level 1 \\
Population D: \> 200 billion people at level 99, 1 billion people at level 2
\end{tabbing}

There are of course features of population D that make it in some respects better than population C, notably that it is more equal and has a worse off person with a higher well-being level. However, there are other features about population D that are much worse than population C, most notably that it involves much less well-being overall. I do not feel like the ARC should give us overwhelming reason to prefer D to C, but this seems to be what it would imply in its standard interpretation. This is because if it is the case that the inequality in A relative to B gives us reason to prefer B to A then it must be the case that the claims of the worse off in A are relevant to those of the better off. However, if this is so then it would seem that the much weaker claims of the better off in C should not be relevant when compared with those of the worse off in C (which are identical to the claims of the worse off in A). If they are not relevant then the fact that there are many more better off people in C compared with A does not matter, because we should simply give absolute priority to the satisfaction of the claims of the worse off people in C.

A further objection to this view arises in the following case:

\begin{tabbing}
Population A:\quad\= \kill
Population E: \> 1 billion people at level 100, 200 billion people at level 1 \\
Population F: \> 10 billion people at level 2, 191 billion people at level 1
\end{tabbing}

If the claims of the worst off in A for choosing B are relevant in comparison with those of the better off in A then, the fact that relevance is directly related to strength according to the ARC means that the claims must also, on this view, be similar in strength, despite the significant difference in the size of benefits involved. This would imply that it should not take the satisfaction of so many claims for us to conclude that it is better to slightly improve the lives of the worse off than greatly diminish the lives of the best off. For instance we might satisfy more strength weighted claims in Population F than in Population E, even though only a small minority of the people in F are better off than they would have been in E, this would still be sufficient to make F better than E, because if these claims are relevant then they must each be `individually' of a similar strength, so that we could conclude that the total strength of claims satisfied in F was greater than that in E. This does not strike me as correct. I believe that when people respond to cases like this they might agree that the inequality in populations like A and E is sufficient to render the otherwise irrelevant claims of the worse off relevant to those of the better off, but they would think that these claims are still, individually at least, much weaker. If we adopted something like my view for cases such as this then we could incorporate this kind of reasoning and say that B is better than A not only because it is more equal, but also because so many people are made better off in B than in A.

Finally, it does appear to be the case that inequality might render an otherwise irrelevant claim (on the grounds of size of benefit alone) relevant, even if it could not render an otherwise relevant claim irrelevant. For instance, let's assume we are considering saving the person with an otherwise very good life from losing their leg or saving somebody with an already terrible life from losing their toe. I think that in this case the fact that the person with a terrible life is already so much worse off might be enough to make their claim to have their toe saved relevant when compared with the claim of the better off person to have their leg saved, even if we would not think that a claim to have a toe saved is usually relevant when compared to having a leg saved. On the other hand, I think that if both individuals needed to have their leg saved then, whilst we may view the claim of the person with the already terrible life as much stronger than that of the person with the otherwise very good life, we would not view any person's claim to have their leg saved as irrelevant to another claim to exactly the same benefit, no matter how much better off they were.

There is no way that I can see for the ARC to come to these kinds of conclusions in its present form. However, if we further modify this view in order to adopt a range conception of relevance that allows inequality, as well as risk, to adjust the size of this range then we can do just this. We can say, in this case, that the worse off a person is, the higher the upper bound of the relevancy of their claim. Relevance will increase in proportion with the strength of their claim, so that even small benefits to the worst off may produce relevant claims in comparison to larger benefits to the best off. However, we can also say that claims that would be relevant for equally well-off people would never become irrelevant for people just because they were unequally well off, i.e. that the lower bound for a claims relevance is the same regardless of how badly off that person is relative to others and depends only on the size of the benefit they would receive.

Implementing this modification would require further changing axiom 4 of the Aggregate Relevant Claims view so that it now says a claim's relevance consists of a range whose upper bound is determined solely by the amount of well-being increase a person would receive and the well-being level of that person (but not the probability that their well-being would be increased) and whose lower bound is determined by the amount of well-being increase that person would receive and the probability that their well-being would be increased (but not their well-being level). One claim is relevant to another if their relevance ranges overlap.

This final version of the view would then look something like:

\begin{enumerate}[label=\arabic*.]
\item Each individual whose well-being is at stake has a claim on you to be helped. An individual for whom nothing is at stake does not have a claim.
\item Individuals' claims compete just in case they cannot be jointly satisfied.
\item An individual's claim is stronger:
\begin{enumerate}[label=\alph*.]
\item the more her well-being would be increased by being aided;
\item the more likely her well-being would be increased by being aided and;
\item the worse off she is relative to others with whom her interests conflict.
\end{enumerate}
\item A claim's relevance consists of a range:
\begin{enumerate}[label=\alph*.]
\item whose upper bound is determined by the amount of well-being increase a person would receive from being aided and how well off she is relative to others with whom her interests conflict and
\item whose lower bound is determined by the amount of well-being increase that person would receive, and the probability that their well-being would be increased.
\end{enumerate}
\par a claim is relevant if its relevancy range overlaps with the relevance range of the strongest competing claim.
\item You should choose an alternative that satisfies the greatest sum of strength-weighted, relevant claims.
\end{enumerate}

Note that in order to interpret the ARC view in this way it will probably be necessary also to substitute axiom 3.b ``An individual's claim is stronger:\ldots{} the lower the level of well-being from which this increase would take place'' with its counterpart from an alternative Competing Claims view presented by Voorhoeve and Fleurbaey (2012): iib ``[An individual's] claim to have a given alternative chosen is stronger:\ldots{} the worse off he is relative to others with whom his interests conflict''. I do not believe this substitution will be problematic. Whilst it is true that this will render the strength of one person's claim inseparable from facts about how well off other people are this will be an especially weak and un-troubling form of inseparability as the only lives that can influence the strength of a person's claims are those of people with whom that person's interests conflict.

\section{The Modified ARC and the Realities of Practical Reason}

The view I have outlined appears to have some radical implications for practical reasoning. In particular, it implies that there may be two different ethical paradigms regarding aggregation, which sit at opposite ends of a continuum. Firstly there are cases where people face certain benefits and harms and are roughly equally well off. In such cases, which are often the ones we start with in many ethical thought experiments, the relevancy range for different claims is narrowly defined, making aggregation between competing claims comparatively less acceptable. In this paradigm, our moral theory will appear more anti-aggregationist, closer to contractarianism or even deontology, because it supports a wide range of trumping considerations, such as never killing or torturing as a means to achieve lesser benefits, no matter how many people stand to benefit. On the other hand, there are cases in which we are dealing with risk and inequality. In such cases, which are usually far more realistic, the relevancy range of claims expands greatly, making aggregation more permissible and moving our moral paradigm towards more utilitarian or consequentialist considerations in which multiple claims must be weighed against one another to produce the greatest benefit overall.

This seems to match our common reaction to different kinds of cases. For instance, the public health planner who must determine how to allocate resources across an entire population is dealing with a significant amount of risk and inequality, so is comparatively more justified in making consequentialist type policies like the overall maximization of QALYs or minimization of DALYs across a population. However, the individual doctor treating each patient one at a time has to face real individuals with more deterministic outcomes and probably less inequality as well. Hence, they may be justified in doing less aggregation and having a more strict prioritization of treatments based on who has the greatest need, even if this produces less well-being or health overall. Generally, I think that this modified version of the ARC would justify more consequentialist type policies at public policy / political level, but less consequentialist type decisions in personal and professional ethics. It may therefore give us some basis for starting to mediate between competing normative frameworks at different levels of practical ethics, which I think would be a very good thing indeed.

\section{Conclusion}

Cases involving risk are often taken to complicate questions about distributive fairness, and are thus typically excluded from theories such as Aggregate Relevant Claims. This is undoubtedly the case, however the incorporation of risk and probability into these theories does not require any radical revision. On the other hand, carefully considering the implications of these cases can reveal further possibilities for how these theories should work in practice that may serve to make them more practically useful and shed new light on their relationship with how practical reason mediates between different kinds of ethical thinking. In particular, applying the Aggregate Relevant Claims view to cases involving risk suggests that the relevance of a claim may not be expressible in terms of a single value, but a range that can expand and contract depending on the circumstances in which the theory is being applied. This mechanism will make our moral thinking more varied, indicating which cases are more or less suitable for the aggregationist approach to fairness, and thus potentially mediating real disagreements about how we should respond to risk.

\appendix
% SJ's choice (not APA): the appendix runs straight on instead of starting a new page
\let\apaclearpage\clearpage\let\clearpage\relax
\section{Can We Apply the ARC to Future Generations?}
\let\clearpage\apaclearpage

\hspace*{\parindent}Not all cases of risk involve only actually existing people. Many of the most pressing global challenges involve risk to possible future people as well, either because they will make conditions in the future worse for the people who live then or because they are cases of Existential Risk, in which no future people may come to exist at all. The proponents of Competing Claims views have argued that theories of distributive fairness, like theirs, should not apply to such cases, because ``a person who does not come into existence cannot have been wronged. Only people who exist at some time can be wronged by not having their interests properly promoted.'' (Voorhoeve \& Fleurbaey, 2016, p. 948).

However, not all philosophers accept this claim. For instance, Derek Parfit argued that existence with a good life can plausibly be said to be `more good' for someone than non-existence, constituting an `existential' benefit, even if not a `comparative' one (Parfit, 2017). This argument would place my `interests' in a mediating position between what is good for me and what can `wrong me'. I.e. there are cases, involving merely possible future people, in which we can say that I have done what is less good for a person, but not wronged them, because in order to wrong them it is necessary for me to fail to `properly promote' their interests, and people do not have interests in outcomes in which they do not exist. On this view, in order to say that merely possible people do not have claims, we would need to show that fairness does not concern what is good for people, but only what is `in their interest', so that even if non-existence is less good for someone, it still cannot ground a claim for existence because they have no interests when they don't exist and so have no claim to what is good for them.

Voorhoeve and Fleurbaey develop one such argument. While they do not consider Parfit's view, they develop an alternative of their own based on the notion that existence may be a good for people on the grounds of what they call `Existence Comparativism', namely that ``It is coherent to appeal to an existing person's preferences in order to judge that her existence is better for her than her never existing.'' Hence we could argue that how good a state would be for a person if they came into existence in it could form the basis of a claim from that possible person in that state against other alternatives in which they did not exist (Voorhoeve \& Fleurbaey, 2016, p. 945).

However, they believe that this claim can only make sense when compared with an alternative and that the only way to make sense of this is to hold that non-existence is equivalent to existence at a well-being level of zero. Hence, ``One ought to treat a case in which a possible person doesn't exist in a state of the world just like a case in which this person exists in that state with a well-being level zero.'' (Voorhoeve \& Fleurbaey, 2016, p. 946).

Next, they evaluate this possible view against a case involving two possible outcomes. In the first outcome, which has a probability of p, Jane exists at level 70+d whilst Kathy does not exist. In the second outcome, which has a probability of 1-p, Kathy exists at level 70 and Jane does not exist. The moral choice under consideration is what value of p would be best, where p can be anything between 0 and 1. Their own view of this case is that ``if Jane's quality of life, were she to exist, would be as good as Kathy's would be if she were to exist instead (d = 0), then it is a matter of indifference which probability you choose. Moreover, if Jane would have a better life than Kathy would have (d > 0), then it would be best if you maximized the chance that Jane would come into existence and chose p = 1.'' (Voorhoeve \& Fleurbaey, 2016, pp. 946--947)

However, Voorhoeve and Fleurbaey argue that this conclusion is incompatible with most theories of distributive fairness if we view possible people as having claims to exist, because these views will be misled by the ex-ante value of these two prospects to the people involved. The higher d becomes the less a \textit{prioritarian} will care about Jane's well-being relative to Kathy's, motivating them to try more to help Kathy by increasing the probability that she will exist. Similarly, the higher d becomes, the more an \textit{egalitarian} will try to equalize the value of the two outcomes between Jane and Kathy, again by increasing the probability of Kathy's coming into existence.

Given this, they argue that distributive fairness is incompatible with the idea that non-existent people can be wronged. They conclude that ``if one chooses to create Jane for sure (p = 1), then Kathy will never be wronged, because she will never be. This contrasts sharply with a scenario in which Kathy will exist for sure alongside Jane, but either, with probability p, has well-being level 0, or, with probability 1 -- p, has a well-being level of 70.'' (Voorhoeve \& Fleurbaey, 2016, p. 948)

This argument only shows that we should accept neither egalitarianism nor (ex-ante) prioritarianism applied to possible future people. However, Voorhoeve and Fleurbaey have already shown why we should not accept either of these views with regard to present generations either (Voorhoeve \& Fleurbaey, 2012). So, perhaps we might still get the correct results by applying a different theory, such as ARC, instead?

One way to do this is to ignore these theories' egalitarian aspects (i.e. that the strength and relevance of people's claims depends on their well-being level), and we actually have good, independent, reasons to do so and only focus on the strength of possible people's claims in terms of the size of their existential benefits, namely that non-existent people do not have a well-being level to consider. If we then consider all the different values of p we might choose in the case of Kathy and Jane: between any two of these values for p we would have to say that the difference this makes to the value of the final prospect for Jane is greater than the difference it makes for Kathy (so long as d is greater than 0). So ARC would justify the conclusion that Voorhoeve and Fleurbaey want to justify, i.e. that what fairness demands is that we simply make it that Jane is the one who exists if d is greater than 0 and are indifferent to all values of p if d is equal to 0.

It may be that Voorhoeve and Fleurbaey don't consider views like ARC because they view their egalitarian character as central to their appeal, and thus see no reason to ignore it in cases involving future generations. However, I think the real problem is in their statement that ``one ought to treat a case in which a possible person doesn't exist in a state of the world just like a case in which this person exists in that state with a well-being level zero.'' That seems like the right thing to say if we accept existence comparativism. However, as Parfit showed, we can also say that a person's interests are at stake merely because an outcome is more or less good for them, but that these are only a subset of their wider interests, which also extend to whether that outcome is better or worse for them as well. If we say this then there is absolutely no reason why we would have to treat merely possible people `just like' actually existing ones. We could instead say something like `one ought to treat a case in which a possible person doesn't exist in a state of the world \textit{similarly} to a case in which this person exists in that state with a well-being level zero.'

In particular, I think we should go on to say that in both of these cases, these outcomes are less good for this person than one in which they would exist with a good life. However, the situation in which they actually exist with a well-being level zero is also worse for this person than that in which they would have a high quality of life, whilst the situation in which they do not exist, whilst still being less good for them, is not worse. Then we could develop our view in many other ways, for instance by claiming that `what is good for a person is in their interests, and we do not treat a person's interests fairly if we fail to do what is good for them when we have the opportunity to do so.' However, not existing, since it is neither better nor worse for me than existing, does not make me `worse off relative to others who do exist'. Hence, whilst I can have claims based on the size of a potential existential benefit I may receive I cannot have a claim based on non-existence leaving me less well off than other people. To put it another way, whilst I have a claim to be brought into existence, the strength of this claim only depends upon how well off I will be in the outcome in which I exist, and not how I will fare relative to others in that outcome.

So I can see no real objection to applying ARC (or other competing claims views) to future generations. I therefore think these views could have a valuable role to play in our evaluation of pressing global challenges like existential risk.

% back to ordinary APA headings for the reference list
\makeatletter\let\section\old@apaSeven@section\makeatother
\section{References}
\begin{list}{}{\setlength{\leftmargin}{0.5in}\setlength{\itemindent}{-0.5in}\setlength{\itemsep}{0pt}\setlength{\parsep}{0pt}}\linespread{2}\selectfont
\item Lazar, S. (2018). Limited aggregation and risk. \textit{Philosophy \& Public Affairs}, \textit{46}(2), 117--159.
\item Otsuka, M. (2015). Risking life and limb: How to discount harms by their improbability. In I. G. Cohen, N. Daniels, \& N. Eyal (Eds.), \textit{Identified versus statistical lives: An interdisciplinary perspective} (pp. 77--93). Oxford University Press.
\item Parfit, D. (2017). Future people, the non-identity problem, and person-affecting principles. \textit{Philosophy \& Public Affairs}, \textit{45}(2), 118--157.
\item Scanlon, T. M. (1998). \textit{What we owe to each other}. Harvard University Press.
\item Voorhoeve, A. (2014). How should we aggregate competing claims? \textit{Ethics}, \textit{125}(1), 64--87.
\item Voorhoeve, A., \& Fleurbaey, M. (2012). Egalitarianism and the separateness of persons. \textit{Utilitas, 24}(3), 381--398.
\item Voorhoeve, A., \& Fleurbaey, M. (2016). Priority or equality for possible people? \textit{Ethics}, \textit{126}(4), 929--954.
\end{list}

\end{document}
